\documentclass[tikz,border=0pt]{standalone}
\usepackage{tikz}
\usepackage{circuitikz}
\usetikzlibrary{calc, positioning, fit, shapes, arrows.meta, backgrounds}

\begin{document}

% Define colors matching the original figure's visual style
\definecolor{bgGreen}{RGB}{220, 235, 215}
\definecolor{boxOrange}{RGB}{253, 226, 171}
\definecolor{blockBlue}{RGB}{124, 166, 223}
\definecolor{pixelGreen}{RGB}{165, 215, 165}
\definecolor{arrowPurple}{RGB}{112, 48, 160}
\definecolor{busRow}{RGB}{80, 110, 180}
\definecolor{busCol}{RGB}{180, 80, 80}

\definecolor{mBlue}{HTML}{1E88E5}
\definecolor{mLightBlue}{HTML}{BBDEFB}
\definecolor{mDarkBlue}{HTML}{0D47A1}

\definecolor{mGreen}{HTML}{43A047}
\definecolor{mLightGreen}{HTML}{C8E6C9}

\definecolor{mOrange}{HTML}{FB8C00}
\definecolor{mLightOrange}{HTML}{FFE0B2}

\definecolor{mGray}{HTML}{ECEFF1}

\begin{tikzpicture}[
    font=\sffamily,
    block/.style={draw=black, fill=blockBlue, thick, align=center},
    dashedbox/.style={draw=black, dashed, thick, fill=boxOrange},
    line/.style={draw=black, thick, -Latex},
    background rectangle/.style={fill=bgGreen},
    signalLine/.style={thick, -{Latex[length=2mm]}},
    connector/.style={circ, scale=0.25},
    show background rectangle,
    inner frame sep=20pt
]

    % Adjusted circuitikz scaling for perfect 1.4x1.4 fits
    \ctikzset{
        bipoles/thickness=1.5,
        nodes width=0.08,
        diodes/scale=0.4,
        grounds/scale=0.4,
        transistors/scale=0.35
    }

    % ==========================================
    % FLEXIBLE PARAMETERS
    % ==========================================
    \def\Ncols{5}          % Number of columns in pixel array
    \def\Nrows{5}          % Number of rows in pixel array
    \def\dx{1.4}           % Width of a single pixel cell
    \def\dy{1.4}           % Height of a single pixel cell
    \def\gap{0.3}          % Standard gap between blocks
    \def\blkW{1.2}         % Standard width for vertical blocks 
    \def\blkH{0.8}         % Standard height for horizontal blocks 
    \def\sPad{0.5}         % Padding for the outer dashed boxes

    \def\nmosOffset{0.95}

    % Derived Dimensions for the Grid
    \pgfmathsetmacro{\Warray}{\Ncols*\dx}
    \pgfmathsetmacro{\Harray}{\Nrows*\dy}

    % Y-Coordinates mapping relative to the grid
    \pgfmathsetmacro{\colAmpsY}{-\gap} 
    \pgfmathsetmacro{\colMuxY}{-\gap - \blkH - \gap} 
    \pgfmathsetmacro{\sensorBottom}{\colMuxY - \blkH}

    % X-Coordinates mapping relative to the grid
    \pgfmathsetmacro{\rowAccX}{-\gap}
    \pgfmathsetmacro{\rowDrvX}{\rowAccX - \blkW - \gap}
    
    \def\clockW{3.2} % Width for Clock & Timing / Oscillator
    \pgfmathsetmacro{\clockX}{\rowDrvX - \blkW - \gap}
    
    \def\biasW{1.8} % Width for Bias Generation
    \pgfmathsetmacro{\biasGenX}{\clockX - \clockW - \gap}
    \pgfmathsetmacro{\sensorLeft}{\biasGenX - \biasW}

    % ==========================================
    % 1. BACKGROUND DASHED BOXES
    % ==========================================
    \draw[dashedbox] (\sensorLeft-\sPad, \sensorBottom-\sPad) rectangle (\Warray+\sPad, \Harray+\sPad);
    \node[anchor=south, align=center, font=\bfseries] at ( {(\sensorLeft+\Warray)/2}, \Harray+\sPad ) 
        {Complementary Metal Oxide Semiconductor\\Image Sensor};

    \def\pcbW{1.2}
    \def\pcbGap{1.5}
    \pgfmathsetmacro{\pcbRight}{\sensorLeft - \sPad - \pcbGap}
    \pgfmathsetmacro{\pcbLeft}{\pcbRight - 2*\pcbW - \gap}
    
    \draw[dashedbox] (\pcbLeft-\sPad, \sensorBottom-\sPad) rectangle (\pcbRight+\sPad, \Harray+\sPad);
    \node[anchor=south, align=center, font=\bfseries] at ( {(\pcbLeft+\pcbRight)/2}, \Harray+\sPad ) 
        {Camera\\(Printed Circuit Board)};

    % ==========================================
    % 2. THE PIXEL ARRAY
    % ==========================================

    % Layer A: Array Background & Cell Borders
    \draw[thick, fill=white, rounded corners=2pt] (0,0) rectangle (\Warray, \Harray);
    
    \foreach \x in {1,...,\Ncols} {
        \foreach \y in {1,...,\Nrows} {
            \pgfmathsetmacro{\px}{(\x-1)*\dx}
            \pgfmathsetmacro{\py}{(\y-1)*\dy}
            
            \ifnum\x=4
                \ifnum\y=3
                    \draw[fill=mLightGreen, rounded corners=1pt] (\px+0.05, \py+0.05) rectangle (\px+\dx-0.05, \py+\dy-0.05);
                \else
                    \draw (\px, \py) rectangle (\px+\dx, \py+\dy);
                \fi
            \else
                \draw(\px, \py) rectangle (\px+\dx, \py+\dy);
            \fi
        }
    }

    % Layer B: Modern Bus Routing
    \foreach \y in {1,...,\Nrows} {
        \pgfmathsetmacro{\py}{(\y-1)*\dy}
        % Reset Bus Line (Shifted to clear transistor bodies)
        \draw[busRow] (\rowAccX, \py+1.1) -- (\Warray-0.1, \py+1.1);
        % Select Bus Line
        \draw[busRow] (\rowAccX, \py+0.15) -- (\Warray-0.1, \py+0.15);
    }
    \foreach \x in {1,...,\Ncols} {
        \pgfmathsetmacro{\px}{(\x-1)*\dx}
        \draw[busCol] (\px+1.2, \colMuxY) -- (\px+1.2, \Harray-0.1);
    }

    % Layer C: Circuit Components (Strict Anchor Binding)
    \foreach \x in {1,...,\Ncols} {
        \foreach \y in {1,...,\Nrows} {
            \pgfmathsetmacro{\px}{(\x-1)*\dx}
            \pgfmathsetmacro{\py}{(\y-1)*\dy}
            
            % 1. Component Instantiation
            \node[nmos] (RST_\x_\y) at (\px+0.55, \py+0.95) {};
            \node[nmos] (SF_\x_\y)  at (\px+\nmosOffset, \py+0.65) {};
            \node[nmos] (SEL_\x_\y) at (\px+\nmosOffset, \py+0.3) {};

            % 2. Photodiode & Ground
            \draw (\px+0.3, \py+0.25) node[ground, line width=0.1]{} 
                to[full photodiode] (\px+0.3, \py+0.65) coordinate (PD_\x_\y);

            % 3. Floating Diffusion Node (FD)
            \draw (PD_\x_\y) -- (SF_\x_\y.G);
            \draw (RST_\x_\y.S) -- (\px+0.55, \py+0.65) node[connector]{};

            % 4. VDD Power Delivery Network
            \draw (\px+\nmosOffset, \py+1.3) node[vcc, scale=0.5] (VDD_\x_\y) {};
            \draw (SF_\x_\y.D) -- (VDD_\x_\y);
            \draw (RST_\x_\y.D) |- (\px+\nmosOffset, \py+1.25) node[connector]{};

            % 5. Amplifier Interconnects
            \draw (SF_\x_\y.S) -- (SEL_\x_\y.D);

            % 6. Gate Tap Routing (to Buses)
            \draw (RST_\x_\y.G) -- (\px+0.15, \py+0.95) -- (\px+0.15, \py+1.1) node[connector]{};
            \draw (SEL_\x_\y.G) -- (\px+0.5, \py+0.3) -- (\px+0.5, \py+0.15) node[connector]{};

            % 7. Readout Trace
            \draw (SEL_\x_\y.S) |- (\px+1.2, \py+0.05) node[connector]{};
        }
    }

    % ==========================================
    % 3. INTERNAL CMOS SENSOR BLOCKS
    % ==========================================
    \draw[block] (0, \colAmpsY) rectangle (\Warray, \colAmpsY-\blkH) node[midway] {Column Amps};
    \draw[block] (0, \colMuxY) rectangle (\Warray, \sensorBottom) node[midway] {Column Mux};

    \draw[block] (\rowAccX, 0) rectangle (\rowAccX-\blkW, \Harray) node[midway, rotate=90] {Row Access};
    \draw[block] (\rowDrvX, 0) rectangle (\rowDrvX-\blkW, \Harray) node[midway, rotate=90] {Row Drivers};

    \draw[block] (\clockX, \Harray) rectangle (\clockX-\clockW, \Harray*0.55) node[midway] {Clock \& Timing\\Generation};
    \draw[block] (\clockX, 0) rectangle (\clockX-\clockW, \Harray*0.45) node[midway] {Oscillator};

    \draw[block] (\biasGenX, 0) rectangle (\biasGenX-\biasW, \Harray) node[midway, rotate=90] {Bias Generation};

    \draw[block] (\rowDrvX-\blkW, \colMuxY) rectangle (\rowAccX, \sensorBottom) node[midway] {Gain};
    
    \pgfmathsetmacro{\adcCenter}{ (\biasGenX + \rowDrvX-\blkW)/2 }
    \def\adcW{1.2}
    \draw[block] (\adcCenter-0.5*\adcW, \colMuxY) rectangle (\adcCenter+0.5*\adcW, \sensorBottom) coordinate[midway] (ADC_mid);
    
    \draw[block] (\biasGenX-\biasW, \colMuxY) rectangle (\biasGenX, \sensorBottom) node[midway] {Line\\Driver};

    % ==========================================
    % 4. CAMERA PCB BLOCKS
    % ==========================================
    \draw[block] (\pcbRight, \sensorBottom) rectangle (\pcbRight-\pcbW, \Harray) node[midway, rotate=90] {Bias Decoupling};
    \draw[block] (\pcbLeft+\pcbW, \sensorBottom) rectangle (\pcbLeft, \Harray) node[midway, rotate=90] {Connector};

    % ==========================================
    % 5. SIGNAL LINES & DATA FLOW ARROWS
    % ==========================================
    \draw[line] (\pcbRight, \Harray/2) -- (\sensorLeft, \Harray/2);

    \draw[line] (0, \colMuxY-0.5*\blkH) -- (\rowAccX, \colMuxY-0.5*\blkH);
    \draw[line] (\rowDrvX-\blkW, \colMuxY-0.5*\blkH) -- (\adcCenter+0.5*\adcW, \colMuxY-0.5*\blkH);
    \draw[line] (\adcCenter-0.5*\adcW, \colMuxY-0.5*\blkH) -- (\biasGenX, \colMuxY-0.5*\blkH);

    \draw[line] (\biasGenX-0.5*\biasW, \sensorBottom) -- ++(0, -1.2) node[below, align=center] {To Frame\\Grabber};

    \node[align=center] (ADC_text) at (\adcCenter, \sensorBottom - \sPad - 1.2) {Analog-to-Digital\\Conversion};
    \draw[-Latex, thick] (ADC_text) -- (\adcCenter, \sensorBottom);

    % ==========================================
    % 6. ANNOTATIONS
    % ==========================================
    \pgfmathsetmacro{\annoX}{\Warray + 2.0}
    
    \node[rotate=90, align=center] (Photon) at (\annoX, \Harray*0.75) {Photon-to-Electron\\Conversion};
    \node[rotate=90, align=center] (Electron) at (\annoX, \Harray*0.25) {Electron-to-Voltage\\Conversion};

    \pgfmathsetmacro{\tarX}{(4-1)*\dx + 0.5*\dx}
    \pgfmathsetmacro{\tarY}{(3-1)*\dy + 0.5*\dy}
    \coordinate (GreenPix) at (\tarX, \tarY);

    \draw[-Latex, thick, arrowPurple] (Photon.north) -- (GreenPix);
    \draw[-Latex, thick, arrowPurple] (Electron.north) -- (GreenPix);

\end{tikzpicture}
\end{document}